
\section{\texorpdfstring{$A_\infty$}{A-infinity} 双模的同伦等价：代数框架}
\label{sec:5-2}

\subsection{弱滤过 \texorpdfstring{$A_\infty$}{A-infinity} 双模同态}
\label{subsec:5-2-1}

第 3.7.1、3.7.2 小节定义了
$\Lambda_{0,\mathrm{nov}}$ 上的滤过 $A_\infty$ 双模
及滤过 $A_\infty$ 双模同态的基本概念。
后续应用需要同时考虑
$\Lambda_{0,\mathrm{nov}}$
与 $\Lambda_{\mathrm{nov}}$ 两种系数环，
因此先讨论 $\Lambda_{\mathrm{nov}}$ 上的
$A_\infty$ 双模同态。

设 $(C_i,\mathfrak m)$ 与
$(C'_i,\mathfrak m')$（$i=0,1$）
是 $\Lambda_{0,\mathrm{nov}}$ 上的
滤过 $A_\infty$ 代数，
$\mathfrak f^{(i)}:C_i\to C'_i$
是它们之间的滤过 $A_\infty$ 同态。
再设 $(D,\mathfrak n)$ 与
$(D',\mathfrak n')$ 分别是
$(C_1,C_0)$ 与 $(C'_1,C'_0)$ 上
以 $\Lambda_{0,\mathrm{nov}}$ 为系数的
滤过 $A_\infty$ 双模。
置
\[
\widetilde D
=D\otimes_{\Lambda_{0,\mathrm{nov}}}
\Lambda_{\mathrm{nov}},
\]
这是 $\Lambda_{\mathrm{nov}}$ 上的
滤过 $A_\infty$ 双模。
注意，$\widetilde D$ 对每个
$\lambda\in\mathbb R$ 都有能量滤过
$F^\lambda\widetilde D$，
而 $D$ 的能量滤过 $F^\lambda D$
只在 $\lambda\geq0$ 时定义。
以同样方式定义 $\widetilde D'$。

\begin{definition}{5.2.1}\label{def:5-2-1}
设 $\widetilde D,\widetilde D'$ 是上述
$\Lambda_{\mathrm{nov}}$ 上的滤过 $A_\infty$ 双模。
在 $(\mathfrak f^{(1)},\mathfrak f^{(0)})$ 之上的
\emph{弱滤过 $A_\infty$ 双模同态}
\liftIndex{I378}
$\widetilde D\to\widetilde D'$，
是满足下述性质的一族
$\Lambda_{\mathrm{nov}}$-模同态
\[
\begin{aligned}
\varphi_{k_1,k_0}:{}
&B_{k_1}(C_1[1])
 \widehat\otimes\widetilde D[1]
 \widehat\otimes B_{k_0}(C_0[1])\\
&\longrightarrow\widetilde D'[1].
\end{aligned}
\]
这里的张量积均取在
$\Lambda_{0,\mathrm{nov}}$ 上。

\par\noindent\textup{(5.2.2.1)}\quad
存在不依赖于 $k_1,k_0$ 的常数 $c\geq0$，
使得
\[
\begin{aligned}
\varphi_{k_1,k_0}\bigl(&
F^{\lambda_1}B_{k_1}(C_1[1])
\widehat\otimes F^\lambda\widetilde D[1]
\widehat\otimes F^{\lambda_0}B_{k_0}(C_0[1])
\bigr)\\
&\subset F^{\lambda_1+\lambda+\lambda_0-c}
\widetilde D'[1].
\end{aligned}
\]

\par\noindent\textup{(5.2.2.2)}\quad
由 $\varphi_{k_1,k_0}$ 与 $\mathfrak f^{(i)}$
诱导的余模同态（见定义 3.7.7）
\[
\begin{aligned}
\widehat\varphi
=\sum_{k_1,k_0}\widehat\varphi_{k_1,k_0}:{}
&\widehat B(C_1[1])
 \widehat\otimes\widetilde D[1]
 \widehat\otimes\widehat B(C_0[1])\\
&\longrightarrow
\widehat B(C'_1[1])
 \widehat\otimes\widetilde D'[1]
 \widehat\otimes\widehat B(C'_0[1])
\end{aligned}
\]
是上链映射，即
$\widehat\varphi\circ\widehat d
=\widehat d'\circ\widehat\varphi$；
$\widehat d,\widehat d'$ 如式 (3.7.2) 所定义。
我们把这一双模同态简记为
$\varphi:\widetilde D\to\widetilde D'$。
\end{definition}

若须指明条件 (5.2.2.1) 中的常数 $c$，
就称 $\varphi=\{\varphi_{k_1,k_0}\}$
为 $c$-弱滤过 $A_\infty$ 双模同态\liftIndex{I319}，
并称 $c$ 为\emph{能量损失}\liftIndex{I183}。

设 $(C_i,\mathfrak m)$、$(C'_i,\mathfrak m')$
带有幺元，且 $\mathfrak f^{(i)}$、
$(D,\mathfrak n)$、$(D',\mathfrak n')$
均含幺（见定义 3.7.5）。
与第 3.7.2 小节一样，
若弱滤过双模同态
$\varphi:\widetilde D\to\widetilde D'$
满足式 (3.7.11)，便称它\emph{含幺}。
此外，若两对代数都带有同伦幺元，
且 $\mathfrak f^{(i)}$、
$(D,\mathfrak n)$、$(D',\mathfrak n')$
均同伦含幺（见定义 3.7.5），
则当 $\varphi$ 可扩张为
在 $(\mathfrak f^{(1),+},\mathfrak f^{(0),+})$
之上的含幺弱滤过 $A_\infty$ 双模同态时，
称 $\varphi$\emph{同伦含幺}。
参见定义 3.7.12。

为了给弱滤过 $A_\infty$ 双模同态
规定有隙条件，
先引入条件 3.1.6 中幺半群
$G\subset\mathbb R_{\geq0}\times2\mathbb Z$
的相应概念。

\begin{definition}{5.2.3}\label{def:5-2-3}
设 $G\subset\mathbb R_{\geq0}\times2\mathbb Z$
是满足条件 3.1.6 的幺半群。
若 $G'\subset\mathbb R\times2\mathbb Z$
满足：任取 $g\in G$ 与 $g'\in G'$，
均有 $g+g'\in G'$，
就称 $G'$ 为\emph{$G$-集合}\liftIndex{I182}。
\end{definition}

下文考虑满足以下条件的 $G$-集合 $G'$：
\par\noindent\textup{(5.2.4.1)}\quad
存在 $c\geq0$，使任意
$(\lambda,\mu)\in G'$ 都满足 $-c\leq\lambda$。
\par\noindent\textup{(5.2.4.2)}\quad
令 $\pi:\mathbb R\times2\mathbb Z\to\mathbb R$
为投影，则 $\pi(G')\subset\mathbb R$ 是离散集。
\par\noindent\textup{(5.2.4.3)}\quad
对任意 $\lambda\in\mathbb R$，
$G'\cap(\{\lambda\}\times2\mathbb Z)$
是有限集。

\begin{definition}{5.2.5}\label{def:5-2-5}
设 $\varphi:\widetilde D\to\widetilde D'$
是位于 $\mathfrak f^{(i)}:C_i\to C'_i$
（$i=0,1$）之上的弱滤过 $A_\infty$ 双模同态，
并假定 $\mathfrak f^{(i)}$ 为 $G$-有隙同态。
取满足 (5.2.4.1)--(5.2.4.3) 的
$G$-集合 $G'$。
在 $G'$ 为无限集时，
可将其元素枚举为
$G'=\{(\lambda'_i,\mu'_i)\mid
i=0,1,2,\ldots\}$，
其中 $\lambda'_i\uparrow\infty$；
有限集则以下采用有限和\translatornote{原书实体页 285
把任意满足 (5.2.4) 的 $G'$ 一律
写成能量趋于无穷的无限序列，
但条件 3.1.6 允许 $G=\{(0,0)\}$，
从而也允许有限的 $G'$。
这里补出有限情形，见疑误 A-0268。}。
若定义 5.2.1 中的分量同态
$\varphi_{k_1,k_0}$
可分解为\translatornote{原书实体页 285
在此误引 (5.2.2.1)，
该编号属于能量滤过条件而非分量同态的源靶式；
见疑误 A-0267。}
\[
\varphi_{k_1,k_0}
=\sum_i T^{\lambda'_i}e^{\mu'_i/2}
\varphi_{k_1,k_0,i},
\]
其中
\[
\begin{aligned}
\varphi_{k_1,k_0,i}:{}
&B_{k_1}(\overline C_1[1])
\otimes\overline D[1]
\otimes B_{k_0}(\overline C_0[1])\\
&\longrightarrow\overline D'[1]
\end{aligned}
\]
都是 $R$-模同态，
便称 $\varphi$ 为\emph{$G'$-有隙}\liftIndex{I120}。
这里把
$\widetilde D=\overline D\otimes_R
\Lambda_{\mathrm{nov}}$ 作同一识别。
\end{definition}

为弱滤过 $A_\infty$ 双模同态
定义有隙条件时使用 $G$-集合；
其他情形的有隙条件则使用幺半群 $G$。
